% Source: https://artofproblemsolving.com/community/c3h3300433_i_found_the_exact_value_of_sin_of_1_degree
% Stable identity: aops-topic-3300433-q01
The exact values for $18^\circ$ and $15^\circ$ show that $c=\sin18^\circ\cos15^\circ-\cos18^\circ\sin15^\circ=\sin3^\circ$. The triple-angle identity gives $3\sin1^\circ-4\sin^31^\circ=c$, so $\sin1^\circ$ is a root in the stated interval. The function $4s^3-3s+c$ has derivative $12s^2-3<0$ there; hence at most one root. It changes from $c>0$ at zero to $c-1<0$ at $\dfrac{1}{2}$, proving existence as well. This supplies an exact algebraic characterization without ambiguous complex cube-root choices.
\begin{asy}
usepackage("amsmath");
size(12cm);
pair O=(0,0),A=dir(20),B=dir(60);draw(circle(O,1),gray);draw((0,0)--(1.2,0),Arrow);draw(O--A^^O--B,blue);draw(A--(A.x,0)^^B--(B.x,0),dashed);draw(arc(O,.25,0,20),blue);draw(arc(O,.5,0,60),red);label("$\theta$",.32dir(10),SE);label("$3\theta$",.6dir(40),NW);label("$\sin\theta$",(A.x,A.y/2),E);label("$\sin3\theta$",(B.x,B.y/2),W);label("$O$",O,SW);label("Generic angle diagram; specialize to $\theta=1^\circ$",(0,-1.3));label("$\sin3\theta=3\sin\theta-4\sin^3\theta$",(0,-1.65));
\end{asy}
